\chapter{S-expressions and rewrite rules}\label{ch:7}
\section*{The working problem}
Choose a representation that can round-trip arbitrary atoms, then make transformation policy explicit.

SExpr distinguishes atoms from lists and tracks quoted atoms. RewriteRule operates on AST values. TermRewriteSystem and EGraph, covered later, attach different traversal and equivalence semantics to superficially similar rule notation.

\section*{Three distinct rewrite engines}
RewriteSet traverses AST values bottom-up and repeats passes. A matching rule marks the pass changed even if its replacement is structurally equal, so a rule that always matches can consume max\_iterations without making progress. The returned tree does not include a stop report. Compare structural equality yourself or use the explicit term reducer when the application must distinguish a normal form from a limit stop.

\section*{In this chapter}
\begin{enumerate}
\item Parsing one S-expression
\item Quoted atom round trips
\item Several top-level forms
\item AST conversion
\item Validated rule notation
\end{enumerate}
Read the workshops in order, then use the chapter review to connect the small examples to an application boundary. Each example is independently runnable; it does not depend on variables introduced in an earlier listing.

\clearpage
\workshop{7.1}{Parsing one S-expression}{C++20}
The singular parser requires exactly one expression. This is a useful boundary for a configuration value or one rewrite declaration. Use the plural parser for a file containing several top-level forms instead of manually splitting on newlines.

\lstinputlisting[language=C++,caption={Example ex07\_01: C++}]{examples/cpp/ex07_01.inc}
\section*{Execution, step by step}
The head add is an ordinary first atom in the list. Operator meaning is supplied by the consuming language.

\section*{Read the implementation}
The relevant source entry is parse\_sexpr in vendor/DSLtk.hpp, around line 9419. Follow the public call into its result construction or dispatch branch. Identify the input ownership, the success condition and the exact failure path before changing this example.

The assertions are the acceptance contract for this workshop. Run the complete program, then change one input at a time. When a change fails, first identify whether the failure belongs to the native build, structural recognition, language checking, execution or the domain result.

\clearpage
\workshop{7.1}{Parsing one S-expression}{D}
D uses module imports and exclamation-mark template instantiation where C++ uses includes, namespaces and angle brackets. The runnable file includes the required imports; the listing shows only its main body.

\lstinputlisting[language=D,caption={Example ex07\_01: D}]{examples/d/ex07_01.inc}
\section*{Failure and design boundary}
S-expression parsing validates structure and escapes, not arity or the semantic validity of an operator.

\section*{Try it yourself}
Reject add with only one operand in a separate semantic validation pass.

\section*{Review the result}
Keep the original accepted fixture and add the new boundary fixture beside it. Compare observable values and diagnostic categories across the two hosts. Do not label the D run as passed until a D compiler has executed it; the supplied edition records source review separately from runtime evidence.

\clearpage
\workshop{7.2}{Quoted atom round trips}{C++20}
Quoted atoms preserve strings that contain whitespace or delimiters. SExpr.atom chooses quoting when necessary, and dump escapes special characters. This makes it a better interchange building block than raw concatenation of AST dump strings.

\lstinputlisting[language=C++,caption={Example ex07\_02: C++}]{examples/cpp/ex07_02.inc}
\section*{Execution, step by step}
The serialized form includes quotes, while the parsed text contains the original payload without those delimiters.

\section*{Read the implementation}
The relevant source entry is SExpr in vendor/DSLtk.hpp, around line 9236. Follow the public call into its result construction or dispatch branch. Identify the input ownership, the success condition and the exact failure path before changing this example.

The assertions are the acceptance contract for this workshop. Run the complete program, then change one input at a time. When a change fails, first identify whether the failure belongs to the native build, structural recognition, language checking, execution or the domain result.

\clearpage
\workshop{7.2}{Quoted atom round trips}{D}
D uses module imports and exclamation-mark template instantiation where C++ uses includes, namespaces and angle brackets. The runnable file includes the required imports; the listing shows only its main body.

\lstinputlisting[language=D,caption={Example ex07\_02: D}]{examples/d/ex07_02.inc}
\section*{Failure and design boundary}
Quoted status can affect pattern-variable interpretation. A quoted question-mark atom should remain literal data.

\section*{Try it yourself}
Round-trip a string containing a quote, a backslash and a newline.

\section*{Review the result}
Keep the original accepted fixture and add the new boundary fixture beside it. Compare observable values and diagnostic categories across the two hosts. Do not label the D run as passed until a D compiler has executed it; the supplied edition records source review separately from runtime evidence.

\clearpage
\workshop{7.3}{Several top-level forms}{C++20}
A rule file naturally contains multiple expressions. parse\_sexpressions performs structural reading across whitespace and comments. It avoids the common mistake of treating a newline inside a quoted string as a record boundary.

\lstinputlisting[language=C++,caption={Example ex07\_03: C++}]{examples/cpp/ex07_03.inc}
\section*{Execution, step by step}
Two independent lists become two array entries. Their separation does not require one form per physical line.

\section*{Read the implementation}
The relevant source entry is parse\_sexpressions in vendor/DSLtk.hpp, around line 9406. Follow the public call into its result construction or dispatch branch. Identify the input ownership, the success condition and the exact failure path before changing this example.

The assertions are the acceptance contract for this workshop. Run the complete program, then change one input at a time. When a change fails, first identify whether the failure belongs to the native build, structural recognition, language checking, execution or the domain result.

\clearpage
\workshop{7.3}{Several top-level forms}{D}
D uses module imports and exclamation-mark template instantiation where C++ uses includes, namespaces and angle brackets. The runnable file includes the required imports; the listing shows only its main body.

\lstinputlisting[language=D,caption={Example ex07\_03: D}]{examples/d/ex07_03.inc}
\section*{Failure and design boundary}
Catch structural reader exceptions at the file boundary and attach a filename. Do not report every malformed file as an empty rule set.

\section*{Try it yourself}
Place comments between forms and check that the list count remains unchanged.

\section*{Review the result}
Keep the original accepted fixture and add the new boundary fixture beside it. Compare observable values and diagnostic categories across the two hosts. Do not label the D run as passed until a D compiler has executed it; the supplied edition records source review separately from runtime evidence.

\clearpage
\workshop{7.4}{AST conversion}{C++20}
The conversion layer gives a conventional list representation to tagged AST nodes. A leaf becomes a tag followed by an atom; an inner node becomes a tag followed by converted children. Preserve the distinction when designing schemas.

\lstinputlisting[language=C++,caption={Example ex07\_04: C++}]{examples/cpp/ex07_04.inc}
\section*{Execution, step by step}
The payload remains one quoted atom rather than two child nodes. This is the key reason to convert structurally before rendering.

\section*{Read the implementation}
The relevant source entry is ast\_to\_sexpr in vendor/DSLtk.hpp, around line 9427. Follow the public call into its result construction or dispatch branch. Identify the input ownership, the success condition and the exact failure path before changing this example.

The assertions are the acceptance contract for this workshop. Run the complete program, then change one input at a time. When a change fails, first identify whether the failure belongs to the native build, structural recognition, language checking, execution or the domain result.

\clearpage
\workshop{7.4}{AST conversion}{D}
D accepts string template arguments directly. ASTNode remains a value with explicit deep-copy behavior, so changing a copied tree must not mutate the original. Keep tests for that behavior when upgrading the module.

\lstinputlisting[language=D,caption={Example ex07\_04: D}]{examples/d/ex07_04.inc}
\section*{Failure and design boundary}
Not every arbitrary S-expression has a valid tagged AST interpretation. The converter expects a literal atom head for a list.

\section*{Try it yourself}
Try a list whose first element is another list and handle the conversion error explicitly.

\section*{Review the result}
Keep the original accepted fixture and add the new boundary fixture beside it. Compare observable values and diagnostic categories across the two hosts. Do not label the D run as passed until a D compiler has executed it; the supplied edition records source review separately from runtime evidence.

\clearpage
\workshop{7.5}{Validated rule notation}{C++20}
A parsed rewrite declaration records its name, left side and right side. Variables on the right must be bound on the left. This catches a common error before any matching begins and gives a clear separation between reading rules and executing them.

\lstinputlisting[language=C++,caption={Example ex07\_05: C++}]{examples/cpp/ex07_05.inc}
\section*{Execution, step by step}
The rule is a data object. No expression has been transformed merely by parsing the declaration.

\section*{Read the implementation}
The relevant source entry is sexpr\_rule in vendor/DSLtk.hpp, around line 9489. Follow the public call into its result construction or dispatch branch. Identify the input ownership, the success condition and the exact failure path before changing this example.

The assertions are the acceptance contract for this workshop. Run the complete program, then change one input at a time. When a change fails, first identify whether the failure belongs to the native build, structural recognition, language checking, execution or the domain result.

\clearpage
\workshop{7.5}{Validated rule notation}{D}
D uses module imports and exclamation-mark template instantiation where C++ uses includes, namespaces and angle brackets. The runnable file includes the required imports; the listing shows only its main body.

\lstinputlisting[language=D,caption={Example ex07\_05: D}]{examples/d/ex07_05.inc}
\section*{Failure and design boundary}
A structurally valid rule can still be mathematically unsound for your language. Effects, overflow and exceptional values belong in the semantic argument.

\section*{Try it yourself}
Change the right side to ?y and verify that construction rejects the unbound variable.

\section*{Review the result}
Keep the original accepted fixture and add the new boundary fixture beside it. Compare observable values and diagnostic categories across the two hosts. Do not label the D run as passed until a D compiler has executed it; the supplied edition records source review separately from runtime evidence.

\clearpage
\section*{Chapter review: Three distinct rewrite engines}
TermRewriteSystem applies first-order rules in a chosen outermost or innermost strategy. EGraph instead inserts alternatives and merges equivalence classes. These engines share concepts but not operational semantics. Do not present a RewriteSet fixed-point loop as equality saturation. Preserve a written rule contract: required shape, semantic preconditions, allowed effects, integer model and termination measure. Quoting matters because an atom beginning with a question mark can be a variable only when the representation marks it as such.

\section*{Questions to answer in your implementation}
\begin{enumerate}
\item Reject add with only one operand in a separate semantic validation pass.
\item Round-trip a string containing a quote, a backslash and a newline.
\item Place comments between forms and check that the list count remains unchanged.
\item Try a list whose first element is another list and handle the conversion error explicitly.
\item Change the right side to ?y and verify that construction rejects the unbound variable.
\end{enumerate}
\section*{A practical completion test}
Choose one workshop and replace its fixed input with an application boundary: a command-line argument, a file, a configuration field or a typed callback. Write down the accepted input shape and the error outcome. Keep parsing and validation separate from publishing state. Preserve the smallest failing input when a regression is found.

